\section*{第 \ref{ch:b1:enzymes} 章——酶与生化催化}

\begin{solution}{exo:b1:enzymes:1}
它降低\omterm{prop:b1:enzymes:whatitdoes}{活化能}（即\omterm{prop:b1:enzymes:whatitdoes}{过渡态}的高度），
从而同时加快两个方向的速率；它不改变 $\Delta G$、
平衡常数和平衡位置。在
能量曲线上，峰值被降低，两端则不变。
\end{solution}

\begin{solution}{exo:b1:enzymes:2}
$K_m$：速率为 $V_{\max}$ 一半时的\omterm{def:b1:enzymes:enzyme}{底物}浓度（\unit{mol/L}）。
$V_{\max}$：\omterm{def:b1:enzymes:enzyme}{底物}饱和时的速率（\unit{mol/s}，或 \unit{\micro
mol/min}）。$k_{\text{cat}}$：饱和时，每个酶分子每秒
转化的\omterm{def:b1:enzymes:enzyme}{底物}分子数（\unit{s^{-1}}）。$k_{\text{cat}}/K_m$：效率，
即低\omterm{def:b1:enzymes:enzyme}{底物}浓度下的二级速率常数，
（\unit{L.mol^{-1}.s^{-1}}）。
\end{solution}

\begin{solution}{exo:b1:enzymes:3}
无\omterm{def:b1:enzymes:inhibitors}{抑制剂}：截距为 1，故 $V_{\max} = 1$；$x$ 截距为 $-0.5$，故
$K_m = 2$。\omterm{def:b1:enzymes:inhibitors}{竞争性抑制}：$V_{\max} = 1$，$x$ 截距为 $-0.25$，$K_m = 4$。
非\omterm{def:b1:enzymes:inhibitors}{竞争性抑制}：截距为 2，$V_{\max} = 0.5$；$x$ 截距为 $-0.5$，
$K_m = 2$。
\end{solution}

\begin{solution}{exo:b1:enzymes:4}
别构（反馈）调节，历时毫秒；\omterm{prop:b1:enzymes:control}{共价修饰}
（磷酸化），历时数秒至数分钟；酶原的蛋白水解激活，
一次性、按需进行；通过基因表达与
降解控制酶的数量，历时数分钟至数小时。
\end{solution}

\begin{solution}{exo:b1:enzymes:5}
$v_0 = 50[S]/(0.1 + [S])$：8.3、25、45.5 和 \qty{49.5}{\micro mol/min}。
由 $0.9 = [S]/(0.1 + [S])$ 得 $[S] = 0.9\,K_m/0.1 = \qty{0.9}{mmol/L}$。
\end{solution}

\begin{solution}{exo:b1:enzymes:6}
$e^{20\,000/2580} = e^{7.75} = 2300$。对于 $10^{10}$：$\Delta E_a =
RT\ln 10^{10} = 2.58\times 23.0 = \qty{59}{kJ/mol}$。
\end{solution}

\begin{solution}{exo:b1:enzymes:7}
$V_{\max} = \qty{0.12}{mol/L}$（每分钟），在 \qty{1}{mL} 中：
\qty{1.2e-4}{mol/min} $= \qty{2.0e-6}{mol/s}$ 的产物；酶量为
\qty{2e-12}{mol}：$k_{\text{cat}} = 2\times 10^{-6}/2\times 10^{-12} =
\qty{1e6}{s^{-1}}$。效率 $10^6/0.012 = \qty{8e7}{L.mol^{-1}.s^{-1}}$，
与扩散极限相差不超过十倍：几乎每一次
与\omterm{def:b1:enzymes:enzyme}{底物}分子的相遇都是有效的。
\end{solution}

\begin{solution}{exo:b1:enzymes:8}
$1 + 2/K_i = 2$：$K_i = \qty{2}{mmol/L}$。对于 \qty{90}{\%}：
$[S] = 9K_m^{\text{app}}$：无\omterm{def:b1:enzymes:inhibitors}{抑制剂}时为 \qty{9}{mmol/L}，
有\omterm{def:b1:enzymes:inhibitors}{抑制剂}时为 \qty{18}{mmol/L}。
\end{solution}

\begin{solution}{exo:b1:enzymes:9}
抑制第一个承诺步骤会使整条途径立即停止，
且不会浪费中间产物，否则它们会积累；
真正起作用的信号是终产物——其丰度恰恰
正是途径已不再需要的信息——而中间产物的
水平并不能说明需求。
\end{solution}

\begin{solution}{exo:b1:enzymes:10}
若蛋白酶在制造它的\omterm{def:b1:cell-unit-of-life:cell}{细胞}内即具活性，便会消化该\omterm{def:b1:cell-unit-of-life:cell}{细胞}；
以无活性的酶原形式合成则无害，直至肠激酶，
肠壁上的一种酶，在十二指肠中把胰蛋白酶原切为胰蛋白酶，
这里正需要消化；胰蛋白酶随即激活其余
酶原（即一个级联反应）。由于痕量胰蛋白酶即可使该级联
过早启动，因此胰腺将一种专一的胰蛋白\omterm{def:b1:enzymes:inhibitors}{酶抑制剂}
装入颗粒中作为保险；这一安排失效即为胰腺炎。
\end{solution}

\begin{solution}{exo:b1:enzymes:11}
无\omterm{def:b1:enzymes:inhibitors}{抑制剂}：$3^3/(3^3 + 3^3) = 0.50$。有\omterm{def:b1:enzymes:inhibitors}{抑制剂}：$27/(166 + 27) = 0.14$，速率
降至其原值的 \qty{28}{\%}。一种米氏酶，其
$K_m$ 由 3 升至 5.5：$3/6 = 0.50$ 变为 $3/8.5 = 0.35$，即
降至原值的 \qty{70}{\%}。\omterm{def:b1:proteins:allostery}{协同性}使同样的曲线移动
在工作点附近的效果超过两倍，
于是更像开关，而非调光器。
\end{solution}

\begin{solution}{exo:b1:enzymes:12}
单纯的催化剂（如铂表面、酸）会加速许多反应，
且不加区分；而酶只加速一种\omterm{def:b1:enzymes:enzyme}{底物}上的一个反应——
其\omterm{def:b1:enzymes:enzyme}{活性位点}的专一性本身就蕴含着关于
\omterm{def:b1:cell-unit-of-life:cell}{细胞}希望进行何种反应的信息。调节又增加一层：别构位点、
磷酸化位点和酶原肽段告诉酶何时、何地
发挥作用，以应对\omterm{def:b1:cell-unit-of-life:cell}{细胞}的状态。其代价是\omterm{def:b1:proteins:peptide}{蛋白质}很大
（一个十余残基的位点，需数百个残基的蛋白来承载），转化速率
慢于最好的无机催化剂，还需要制造、修饰和
降解它的机制；其好处在于，化学过程仅在
何处、何时有用时才发生，这正是代谢的本质。
\end{solution}

\begin{solution}{pb:b1:enzymes:1}
\textbf{1.} $1/[S]$：2、1、0.5、0.2、0.1、0.05。$1/v_0$：5.0、3.0、2.0、
1.4、1.2、1.1。
\textbf{2.} $1/[S]$ 每变化一步，$1/v_0$ 随之成比例变化：斜率
$(5.0 - 1.1)/(2 - 0.05) = 2.0$；截距 $1.0$。
\textbf{3.} $V_{\max} = 1/1.0 = \qty{1.0}{\micro mol/min}$；$K_m =
\text{斜率}\times V_{\max} = \qty{2.0}{mmol/L}$。
\textbf{4.} $1.0\times 5/(2 + 5) = 0.714$：与实测值相符。
\textbf{5.} $\qty{10e-9}{mol/L}\times\qty{1e-3}{L} = \qty{1e-11}{mol}$；
$V_{\max} = \qty{1e-6}{mol/min} = \qty{1.67e-8}{mol/s}$；
$k_{\text{cat}} = 1.67\times 10^{-8}/10^{-11} = \qty{1670}{s^{-1}}$。
\textbf{6.} $1670/(2\times 10^{-3}) = \qty{8.3e5}{L.mol^{-1}.s^{-1}}$：
比扩散极限低一千倍；平均每一千次碰撞
中只有一次有效。
\textbf{7.} $0.95 = [S]/(2 + [S])$：$[S] = 19\,K_m = \qty{38}{mmol/L}$。
\textbf{8.} $1/v_0$：9.0、5.0、3.0、1.8、1.4、1.2。斜率 $(9.0 -
1.2)/1.95 = 4.0$；截距 $1.0$。
\textbf{9.} $V_{\max}^{\text{app}} = \qty{1.0}{\micro mol/min}$
（不变）；$K_m^{\text{app}} = 4.0\times 1.0 = \qty{4.0}{mmol/L}$。
\textbf{10.} $K_m$ 加倍，$V_{\max}$ 不变：\omterm{def:b1:enzymes:inhibitors}{竞争性抑制}。
\textbf{11.} $4 = 2(1 + 1/K_i)$：$K_i = \qty{1.0}{mmol/L}$。
\textbf{12.} 有\omterm{def:b1:enzymes:inhibitors}{抑制剂}时，$v = V_{\max}[S]/(K_m(1 + [I]/K_i) +
[S])$；速率减半要求 $K_m(1 + [I]/K_i) + [S] = 2(K_m +
[S])$，即 $[I] = K_i(K_m + [S])/K_m$：在 \qty{2}{mmol/L} 时，$[I] =
1\times 4/2 = \qty{2}{mmol/L}$；在 \qty{20}{mmol/L} 时，$[I] = 22/2 =
\qty{11}{mmol/L}$。\omterm{def:b1:enzymes:enzyme}{底物}越多，所需的\omterm{def:b1:enzymes:inhibitors}{抑制剂}也越多。
\textbf{13.} 一个与\omterm{def:b1:enzymes:enzyme}{底物}（或其\omterm{prop:b1:enzymes:whatitdoes}{过渡态}）相似的分子
——一种不能被水解的酯类类似物——结合在
\omterm{def:b1:enzymes:enzyme}{活性位点}上。
\textbf{14.} $1/v_0$：10.0、6.0、4.0、2.8、2.4、2.2；斜率 $4.0$，
截距 $2.0$：$V_{\max}^{\text{app}} = \qty{0.5}{\micro mol/min}$，
$K_m^{\text{app}} = 4.0\times 0.5 = \qty{2.0}{mmol/L}$。
\textbf{15.} $V_{\max}$ 减半，$K_m$ 不变：非\omterm{def:b1:enzymes:inhibitors}{竞争性抑制}；
$2 = 1 + 1/K_i$：$K_i = \qty{1.0}{mmol/L}$。
\textbf{16.} J 结合于\omterm{def:b1:enzymes:enzyme}{活性位点}以外的位点，对游离酶和
$ES$ 的亲和力相同：\omterm{def:b1:enzymes:enzyme}{底物}不与之竞争，
在任何 $[S]$ 下，都有一半的酶分子
失活。
\textbf{17.} 加入 J 后速率加倍（酶量加倍后，其中一半仍
有活性）；加入 I 后速率同样加倍（在任何 $[S]$ 下，速率
都与酶量成正比）：加倍酶量
并不能区分这两种\omterm{def:b1:enzymes:inhibitors}{抑制剂}。
\textbf{18.} \qty{25}{\celsius} 时的速率：$V_{\max}[S]/(K_m + [S])$，其中
$V_{\max} = k_{\text{cat}}[E]$：每升，$1670\times 10^{-8} =
\qty{1.67e-5}{mol/L/s}$；$\times 0.5/2.5 = \qty{3.3e-6}{mol/L/s}$；
在 \qty{37}{\celsius} 时，乘以 $\times 2^{1.2} = 2.3$：\qty{7.7e-6}{mol/L/s}；
在 \qty{1e-12}{L} 中：\qty{7.7e-18}{mol/s} $= \num{4.6e6}$ 个分子
每秒。
\textbf{19.} 无需升高：速率（\num{4.6e6}）超过需求
（\num{3e6}）；若需升高，\omterm{def:b1:cell-unit-of-life:cell}{细胞}可提高\omterm{def:b1:enzymes:enzyme}{底物}浓度，或
使酶磷酸化以降低其 $K_m$，或合成更多的酶。
\textbf{20.} $K_m^{\text{app}} = 2(1 + 0.4/0.2) = \qty{6}{mmol/L}$：速率
$\times 0.5/6.5$ 而非 $0.5/2.5$，为原来的 \qty{38}{\%}，
每秒 \num{1.8e6}——产物抑制自身合成，使之
低于需求，产物随之被消耗、水平下降，从而解除对酶的
抑制：一种自我调节的供应方式。
\textbf{21.} 之前：$0.5/2.5 = 0.20$（占 $V_{\max}$ 的比例）；之后：$0.5/0.9 =
0.56$：接近三倍。
\textbf{22.} 接近零：在 \omterm{def:b1:water-small-molecules:ph}{pH} 5 时，\omterm{def:b1:enzymes:enzyme}{活性位点}的组氨酸被质子化，
无法充当碱，位点中的酸性残基
也被中和；酶的电离状态，也许还有其折叠方式，
都不正确。它是为\omterm{def:b1:eukaryotic-cell:organelle}{细胞溶胶}而非\omterm{def:b1:eukaryotic-cell:endomembrane}{溶酶体}设计的。
\textbf{23.} $K_m$ 变化很小（结合主要依靠位点的其余
部分）；$k_{\text{cat}}$ 下降了数个数量级，因为
组氨酸正是活化水或丝氨酸的催化碱。
\textbf{24.} 在 $K_m$ 附近，速率对\omterm{def:b1:enzymes:enzyme}{底物}的变化敏感
（斜率约为 $V_{\max}/2K_m$）；远高于此时，酶已饱和，且
对之不敏感，过量\omterm{def:b1:enzymes:enzyme}{底物}还将造成渗透和
化学负担。在 $K_m$ 附近工作，使途径可由
供应量来调控。
\textbf{25.} $K_m = \qty{2.0}{mmol/L}$，$V_{\max} = \qty{1.0}{\micro mol/min}$
（\qty{10}{nmol/L} 酶），$k_{\text{cat}} = \qty{1670}{s^{-1}}$；I
为\omterm{def:b1:enzymes:inhibitors}{竞争性抑制}，$K_i = \qty{1.0}{mmol/L}$；J 为非\omterm{def:b1:enzymes:inhibitors}{竞争性抑制}，$K_i =
\qty{1.0}{mmol/L}$。
\end{solution}
